\chapter{Background}
\label{chap:1}

\section{Il Model Context Protocol}\label{introduzione-su-mcp}

L'inserimento dei Large Language Models nei flussi di lavoro è stato
storicamente limitato dalla frammentazione delle connessioni alle fonti
di dati esterne. Prima della diffusione del Model Context Protocol (MCP) 
\autocite{hasan2026modelcontextprotocolmcp}, l'implementazione di un agente 
AI richiedeva lo sviluppo di integrazioni su misura per le API di ogni 
singolo servizio a cui il modello doveva accedere. Seguendo questo approccio, 
all'aumentare dei servizi integrati, le applicazioni AI diventavano molto 
complesse e fortemente accoppiate, e di conseguenza anche fragili e difficili 
da mantenere e scalare.

Nel 2023 si è registrato un primo progresso significativo con
l'introduzione da parte di OpenAI del function calling e dei ``ChatGPT
plugins''. Il primo forniva un meccanismo nativo tramite cui i large
language model (LLM) potevano richiamare degli strumenti esterni, i
secondi, invece, davano la possibilità a sviluppatori di terze parti di
rendere disponibili direttamente i loro servizi al modello sotto forma
di plugin.

Nello stesso periodo altri fornitori AI hanno sviluppato soluzioni
analoghe, ed anche vari framework hanno introdotto delle interfacce
standardizzate per l'accesso agli strumenti da parte degli LLM,
facilitando l'integrazione. Tuttavia, anche questa strategia ha
alimentato un ecosistema frammentato, dove i fornitori di servizi si
sono trovati a realizzare molteplici implementazioni ridondanti per
rendere disponibili i propri strumenti nelle varie piattaforme,
dovendosi adattare alle relative specificità.

Nel novembre 2024 Anthropic ha presentato \textbf{MCP}, un protocollo 
open-source che permette alle applicazioni AI di interagire in modo 
standard con i sistemi software per accedere al contesto e alle capacità 
che forniscono.

Con l'avvento di questo nuovo protocollo si assiste ad un cambiamento di
paradigma radicale «\emph{from tool bindings hardcoded per application
toward an interoperable ecosystem of composable, discoverable network
services}» \autocite{hou2025modelcontextprotocolmcp}.

L'adozione di MCP introduce molteplici vantaggi strutturali
nell'integrazione tra agenti AI e sistemi esterni. Innanzitutto, ogni
server MCP con cui l'agente interagisce, indipendentemente dalla sua
funzione, utilizza un unico protocollo di comunicazione. Dunque, una
volta implementato il supporto a MCP, il client può comunicare con
qualsiasi server MCP esistente. Inoltre, esponendo i
loro servizi tramite un server MCP, gli sviluppatori mantengono l'implementazione
agnostica rispetto ai diversi modelli e framework AI.

Un ulteriore vantaggio consiste nel fatto che gli agenti possono
interrogare direttamente i server MCP per individuare le funzioni
disponibili nel momento corrente. Questo processo di scoperta dinamica
disaccoppia gli agenti dalle funzionalità dei server, superando il
precedente modello basato su integrazioni \emph{hardcoded}, e consentendo a
server ed agenti di evolversi in modo indipendente.

In aggiunta, grazie alla bidirezionalità dei canali di comunicazione
client-server, il protocollo permette di strutturare delle interazioni
più elaborate del semplice utilizzo di uno strumento da parte
dell'agente. Infatti, MCP consente ai server, non solo di inviare
notifiche ai client, ma anche di richiedere dati supplementari
all'utente prima di finalizzare un'operazione.

Infine, MCP è progettato con un'attenzione particolare alla sicurezza,
integrando standard consolidati come OAuth 2.1 per la gestione di
autenticazione e autorizzazione. Il protocollo assicura inoltre che le
credenziali dei servizi esterni restino isolate nel server, evitando che
siano esposte all'agente.

MCP ha visto la sua adozione crescere in modo estremamente rapido a
testimoniare l'efficacia dei benefici sopracitati nelle applicazioni
reali. Attualmente, ad esempio, sfruttando dei server MCP un agente può
accedere al Google Calendar del suo utente per visionarne gli impegni o
può eseguire un comando nel terminale, oppure ancora leggere ed inviare
messaggi su un canale Slack.

\subsection{Architettura}\label{architettura}

MCP è un protocollo basato su un’architettura client-host-server che
prevede tre tipi di attori: \textbf{MCP host}, \textbf{MCP client}, e 
\textbf{MCP server} \autocite{ArchitectureoverviewModelContextProtocol-2026-07-28}.

L'attore principale, ossia l'host, è un agente AI che coordina le 
operazioni e richiede accesso al contesto per fornirlo al modello AI
incapsulato al suo interno. Il contesto a cui l'host deve accedere è 
esposto dagli MCP server, dunque l'host ha necessità di comunicare con 
uno o più di essi. Per farlo, l'host non gestisce le connessioni in 
prima persona, ma istanzia per ciascun server un componente interno 
dedicato, chiamato MCP client. Quest'ultimo è un modulo software che si 
occupa unicamente di mantenere il canale di comunicazione bidirezionale 
tra l'host e uno specifico MCP server.

La \zcref{fig:diagramma_architettura} descrive visivamente l'architettura
di MCP.

\begin{figure}[htbp]
    \centering
    \includegraphics[width=0.5\textwidth]{diagrams/architettura.pdf}
    \caption{Descrizione di alto livello dell'architettura di MCP. L'agente AI (host) comunica con uno o più server MCP tramite i rispettivi client MCP.}
    \label{fig:diagramma_architettura}
\end{figure}

\subsubsection{Gli attori}\label{gli-attori}

Innanzitutto, l'host svolge un ruolo di coordinamento, creando le
istanze del client e gestendo tramite esse le connessioni con i server,
e inoltre fa da tramite tra i server e l'agente AI. Infine, l'host si
occupa di applicare le politiche di sicurezza, e di stabilire se
autorizzare o meno determinate azioni, eventualmente chiedendo conferma
all'utente.

I client, invece, fungono da endpoint per la comunicazione tra gli host
e i server. Ogni client stabilisce e mantiene una connessione dedicata
verso uno specifico server, instradando i messaggi dall'host al server e
viceversa. Inoltre, il client gestisce lo scambio di informazioni sulle
capacità con il server, le notifiche e le sottoscrizioni alle risorse.

Infine, il server integra le informazioni e le funzionalità offerte da
uno o più servizi esterni e le espone agli agenti AI, permettendogli di
estendere il contesto a loro disposizione. I server, inoltre, possono
anche replicare ai messaggi dei client richiedendo l'inserimento di
informazioni aggiuntive da parte dell'utente.

I server MCP possono eseguire sia localmente sulla stessa macchina
dell'host, che remotamente.

Un tipico contesto in cui gli attori suddetti interagiscono è un flusso
di lavoro che ha inizio con l'invio di una richiesta da parte
dell'utente verso l'agente AI. In risposta, l'agente compone un prompt
contenente la richiesta e la lista degli strumenti utilizzabili,
ottenuta interagendo con uno o più server MCP di cui l'agente dispone
attraverso i rispettivi client. Una volta inoltrate tali informazioni al
modello AI, quest'ultimo valuta le funzionalità necessarie per
soddisfare la richiesta, ed indica ai client MCP quali strumenti
invocare e con quali parametri. Infine, dopo aver orchestrato i vari
strumenti per raggiungere il suo obiettivo, il modello elabora una
risposta per l'utente in cui lo informa sull'esito delle operazioni
svolte.

Gli editor di codice basati su AI come Visual Studio Code illustrano
bene l'applicazione pratica dell'architettura di MCP. Gli agenti AI
integrati negli editor possono connettersi tramite degli MCP client a
svariati MCP server. Ad esempio, come si può osservare nella 
\zcref{fig:diagramma_architettura_esempio}, Github Copilot in VS Code può
utilizzare un server locale per accedere al FileSystem e creare,
modificare, e cancellare file, o sfruttare il server remoto di Github
per controllare la storia dei commit e le pull request dei repository.

\begin{figure}[htbp]
    \centering
    \includegraphics[width=0.6\textwidth]{diagrams/architettura_esempio_VS_Code.pdf}
    \caption{Descrizione di alto livello dell'architettura di MCP in un contesto reale. L'agente AI integrato in VS Code comunica con due server MCP, uno locale e uno remoto, tramite i rispettivi client MCP.}
    \label{fig:diagramma_architettura_esempio}
\end{figure}

\subsection{Data layer e transport
layer}\label{data-layer-e-transport-layer}

MCP adotta una struttura modulare che separa nettamente le
responsabilità tra due livelli distinti \autocite{SpecificationModelContextProtocol-2026-07-28}. 
Il data layer individua il formato e la semantica dei messaggi, mentre 
il transport layer si occupa esclusivamente di come i messaggi vengono 
incapsulati e consegnati.

Questa separazione garantisce che la semantica del protocollo rimanga
uniforme su qualsiasi canale di comunicazione, permettendo di introdurre
nuovi trasporti senza alterare la logica applicativa.

\subsection{Data layer}\label{data-layer}

Il data layer di MCP definisce il protocollo per la comunicazione tra
client e server, il quale si basa sulla specifica JSON-RPC 2.0 come
formato di scambio dei messaggi.

\subsubsection{JSON-RPC}\label{json-rpc}

JSON-RPC è un protocollo di chiamata di procedura remota che utilizza
JSON (JavaScript Object Notation) come formato per definire le strutture
dati su cui si basa, e non è vincolato ad un meccanismo di trasporto
specifico.

Il protocollo è strutturato in modo molto semplice e si basa
principalmente su due oggetti \textbf{Request} e \textbf{Response}. Il
primo oggetto rappresenta la chiamata remota che viene inviata dal
client al server, e contiene un campo per il nome della procedura da
invocare e uno per gli eventuali parametri da passarle.

L'oggetto Response, invece, rappresenta l'esito della chiamata remota che
il server restituisce al client, e contiene due campi result e error.
Result memorizza il valore prodotto in output dalla procedura invocata,
e viene incluso nella risposta dal server se l'esecuzione del metodo
termina con successo. Error, invece, è un oggetto strutturato che
descrive l'errore verificatosi nella gestione della chiamata remota, ed è
incluso in Response proprio quando il server riscontra un problema nel
processare la richiesta del client.

Per permettere di associare una risposta del server alla relativa
richiesta, JSON-RPC definisce un campo id in Request e Response. Tale
campo deve essere impostato all'interno di una risposta con lo stesso
valore contenuto nell'id della richiesta a cui fa seguito.

Si riporta qui un esempio basilare di scambio Request-Response tratto
dalla specifica JSON-RPC 2.0 in cui un client effettua una chiamata remota
di un metodo subtract che effettua una sottrazione tra due numeri, e il
server gestisce con successo la richiesta:

\begin{itemize}
\item Richiesta del client: 
    \lstinline|{"jsonrpc": "2.0", "method": "subtract", "params": [42, 23], "id": 1}|

\item Risposta del server: 
    \lstinline|{"jsonrpc": "2.0", "result": 19, "id": 1}|
\end{itemize}

JSON-RPC supporta anche l'invio di notifiche, ossia messaggi
unidirezionali, per cui il client che li invia non necessita, e non
riceve, alcuna risposta da parte del server. Le notifiche
(\textbf{Notification}) sono rese dal protocollo semplicemente come
oggetti Request in cui il campo id viene omesso.

\subsubsection{I messaggi}\label{i-messaggi}

MCP definisce tre tipi fondamentali di messaggi che client e server si
possono scambiare: richieste, risposte, e notifiche.

Le \textbf{richieste} sono dei messaggi finalizzati a richiedere
l'esecuzione di un'operazione da parte di un server o di un client MCP.
Esse si basano sugli oggetti Request di JSON-RPC 2.0, e in aggiunta MCP
impone che si includa sempre al loro interno un id non-nullo, e che
questo sia diverso da qualsiasi id utilizzato dal mittente in una
richiesta precedente per cui non ha ancora ricevuto una risposta.

Le \textbf{risposte} sono dei messaggi che un server o un client invia
come riscontro per una richiesta ricevuta e sono basate sugli oggetti
Response di JSON-RPC 2.0.

Esistono due tipi di risposta: le result responses, che sono inviate
quando l'operazione richiesta viene eseguita con successo e che
contengono il campo result, e le error responses, inviate quando viene
riscontrato un errore nella gestione della richiesta e contenenti il
campo error. Inoltre, si nota che MCP aggiunge rispetto a JSON-RPC il
requisito che l'oggetto result contenga un campo resultType.
Quest'ultimo permette ai client di determinare in modo dinamico la
struttura di result, abilitando il polimorfismo del risultato.

Infine, le \textbf{notifiche} sono dei messaggi unidirezionali che i
client possono inviare ai server e viceversa. Sono definite in modo
sostanzialmente identico alle Notification di JSON-RPC 2.0.

\subsubsection{I pattern di messaggi}\label{i-pattern-di-messaggi}

Le diverse funzionalità del protocollo MCP si realizzano attraverso tre
cosiddetti pattern di messaggi, ovvero specifiche tipologie di
interazione client-server. Ciascuno di questi schemi d'interazione viene
strutturato combinando gli elementi costitutivi descritti nel paragrafo
precedente, ossia i tre tipi di messaggio: richiesta, risposta e
notifica.

Un principio fondamentale osservato da tutti i pattern risiede nella
natura asimmetrica delle interazioni, le quali vengono sempre originate
dal client. Quest'ultimo assume il ruolo di iniziatore inviando
richieste e notifiche, a fronte delle quali il server si limita a
fornire un riscontro. Coerentemente con tale impostazione, ai server è
preclusa la possibilità di avviare autonomamente richieste, così come i
client non emettono risposte.

Il primo pattern è chiamato \textbf{Request and Response} e rappresenta
la modalità di interazione più lineare: a una richiesta inoltrata dal
client il server replica con un oggetto di tipo result o error. Nel
corso dell'elaborazione, il server ha la facoltà di trasmettere
notifiche intermedie, come aggiornamenti sullo stato di avanzamento.

Qualora l'esecuzione di un'operazione richieda al server l'acquisizione
di ulteriori dati dal client si adotta il pattern \textbf{Multi
Round-Trip Requests}. In tale scenario, il server risponde alla
richiesta del client con un risultato di tipo \texttt{input\_required},
inducendo il client a reiterare la richiesta integrando dei parametri
aggiuntivi.

Infine, nel pattern \textbf{Subscribe and Notify}, il client invia una
richiesta con il metodo \path{subscriptions/listen} per essere aggiornato dal
server quando si verificano eventi di un certo tipo, ad esempio quando
viene modificata una particolare risorsa. Nel momento in cui il server
riceve tale richiesta apre uno stream di notifiche del tipo richiesto
verso il client.

\subsubsection{Le primitive}\label{le-primitive}

MCP individua una serie di primitive, le quali rappresentano delle
tipologie di funzionalità (operazioni e contesto informativo) che i
server e i client possono offrirsi reciprocamente.

\subsubsection{Primitive dei server}\label{primitive-dei-server}

MCP permette la \textbf{scoperta dinamica} da parte dell'host delle
funzionalità offerte dal server. Ciò significa che un agente AI, mentre
utilizza un server MCP, può rilevare se le capacità che esso espone
cambiano e come.

Questo è possibile grazie al fatto che per ciascuna primitiva è definito
un metodo \texttt{*/list} che elenca tutte le funzioni disponibili per la
specifica categoria, ed un metodo \texttt{*/get} per ottenere una particolare
funzione appartenente alla determinata categoria.

Inoltre, quando per una certa primitiva cambiano le funzionalità esposte dal 
server, è il server stesso a comunicarlo ai client che si sono “iscritti”
a tale primitiva (message pattern Subscribe and Notify), inviando ad essi
una notifica.

I metodi che i server possono mettere a disposizione rientrano in tre
diverse primitive: strumenti, risorse, e prompt.

Gli \textbf{strumenti} (tool) sono funzioni eseguibili che gli agenti AI
possono invocare per compiere azioni, come interagire con il FileSystem
per modificare un file. Tali strumenti vengono richiamati mediante il
metodo \path{tools/call}.

Le \textbf{risorse} (resources) rappresentano sorgenti dati a cui le
applicazioni AI possono accedere in sola lettura per ottenere
informazioni di contesto, come il contenuto di documenti o dei record
estratti da un database. L'accesso a tali risorse avviene tramite il
metodo \path{resources/read}.

I prompt, infine, consistono in template di istruzioni per illustrare
all'agente AI come utilizzare specifici tool o risorse per svolgere
determinati compiti.

\subsubsection{Primitive dei client}\label{primitive-dei-client}

I client mettono a disposizione un'unica primitiva,
l'\textbf{elicitation}, la quale permette ai server di richiedere
informazioni aggiuntive all'utente dell'agente AI. I server effettuano
tali richieste nella modalità individuata dal pattern Multi Round-Trip
Requests ed utilizzando il metodo \path{elicitation/create}.

Si osserva che in precedenza MCP definiva quattro diverse primitive lato
client, ma tre di queste (sampling, roots, e logging) sono state
deprecate nella versione 2026-07-28 della specifica.

Per una trattazione più dettagliata e corredata da esempi delle primitive 
si rimanda al \zcref{server} per quelle lato server e al \zcref{client} 
per quelle lato client.

\subsubsection{La connessione}\label{la-connessione}

Inizialmente, il protocollo MCP si avvaleva di un modello di tipo
\textbf{stateful}, basato su una procedura di handshake iniziale per
l'apertura della sessione. Durante tale fase, client e server
negoziavano la versione e le capacità, permettendo al server di
mantenere tale contesto informativo per l'intero arco della connessione.
Tuttavia, tale configurazione ha mostrato diverse criticità, tra cui la
difficoltà di scalare orizzontalmente il servizio, la scarsa resilienza
in caso di guasti e l'elevata complessità implementativa sia lato client
che lato server.

Dunque, con l'avvento della specifica 2026-07-28, MCP è evoluto in un
protocollo \textbf{stateless}.

In questo nuovo paradigma, i singoli messaggi di richiesta del client
contengono nel campo \texttt{\_meta} tutti i dati necessari per essere elaborati
dal server: nome e versione del client, versione del protocollo e
capacità operative. Il server gestisce ciascuna richiesta in modo
indipendente, senza desumere alcun contesto da interazioni pregresse,
seppur avvenute all'interno del medesimo canale comunicativo. Qualora
l'esecuzione richieda che uno stato persista tra più richieste, il
client è tenuto a inserire un identificatore esplicito in ogni
successiva invocazione.

\subsubsection{Negoziazione delle
capacità}\label{negoziazione-delle-capacituxe0}

Nella nuova versione del protocollo anche il sistema di negoziazione
delle capacità cambia significativamente: il client specifica in ogni
richiesta le proprie funzionalità nell'attributo
\path{_meta.io.modelcontextprotocol/clientCapabilities}, mentre le capacità
del server sono rilevabili attraverso il metodo \path{server/discover}.

\subsection{Transport layer}\label{transport-layer}

Il transport layer di MCP definisce i meccanismi di trasporto
sottostanti al protocollo di comunicazione, anche detti
\textbf{binding}, per mezzo dei quali i messaggi vengono incapsulati e
trasferiti tra il client e il server. Inoltre, i binding determinano
anche come la connessione tra client e server viene inizializzata e
terminata, come vengono trasmessi i metadati, e come il client può
cancellare una richiesta inoltrata al server.

I meccanismi di trasporto operano in modo del tutto indipendente
rispetto al contenuto dei messaggi, garantendo che i pattern
d'interazione conservino la medesima semantica a prescindere dallo
specifico canale di comunicazione adottato.

MCP prevede due canali di trasporto standard: stdio, per la
comunicazione tra un client ed un server locale, e Streamable HTTP, tra
un client e un server remoto.

\subsubsection{Stdio}\label{stdio}

Il meccanismo di trasporto \textbf{Stdio} prevede che il client esegua
il server MCP come un processo figlio, stabilendo un canale di
comunicazione bidirezionale basato sugli stream standard \texttt{stdin} e \texttt{stdout}.
In tale configurazione, i messaggi vengono scambiati in sequenza,
separati da caratteri di nuova riga, garantendo che \texttt{stdout} sia dedicato
esclusivamente alla trasmissione di pacchetti MCP sintatticamente
validi, mentre eventuali log vengono instradati su \texttt{stderr}. Infine, la
terminazione della comunicazione viene innescata dal client attraverso
la chiusura dello stream di input.

\subsubsection{Streamable HTTP}\label{streamable-http}

Nel trasporto tramite \textbf{Streamable HTTP} il server si configura
come un processo autonomo che espone un unico endpoint HTTP. Ogni
interazione avviata dal client viene veicolata tramite il metodo \texttt{HTTP
POST}, a cui il server può rispondere restituendo un singolo oggetto JSON
o, in alternativa, attivando uno stream SSE (Server-Sent Events)
dedicato alla specifica richiesta. Sotto il profilo della sicurezza, i
server sono obbligati a verificare l'header Origin per prevenire
attacchi di DNS rebinding.

\subsection{Server}\label{server}

I server MCP, come precedentemente illustrato, costituiscono i blocchi
fondamentali tramite cui è possibile estendere il contesto a
disposizione degli LLM. Un server espone in modo standard i dati e le
funzionalità di uno o più servizi alle applicazioni AI, e può offrire
tre tipi di primitive: tool, risorse, e prompt. Di seguito si
approfondiscono tali primitive.

\subsubsection{Tool}\label{tool}

I tool sono funzioni eseguibili che permettono al modello AI di
intraprendere azioni, come effettuare chiamate API verso servizi
esterni. Nel caso degli strumenti, il controllo è affidato al modello
che, analizzando le richieste dell'utente, decide autonomamente se, come
e quando invocare un determinato tool tra quelli resi disponibili dal
server. Ogni tool espone uno schema JSON che ne definisce i parametri di
input richiesti e la struttura dell'output atteso. I risultati
restituiti dal server a seguito dell'invocazione di un tool possono
presentarsi sia in formato non strutturato, come testo o immagini, che
in formato strutturato, quindi come dei valori JSON. A causa della
natura attiva dei tool e della loro capacità di alterare lo stato dei
sistemi esterni, il protocollo suggerisce di affiancare a questa
primitiva dei meccanismi di sicurezza di tipo human-in-the-loop, che
delegano all'utente finale la conferma e l'autorizzazione prima di
finalizzare l'esecuzione delle operazioni.

\paragraph{Esempio}\label{esempio-tool}

Al fine di illustrare l'impiego dei tool, si considera un agente AI a
cui l'utente chiede le previsioni meteo per la città di Torino.
L'interazione con il server MCP del meteo necessaria per assolvere a
questo compito esemplifica anche direttamente il message pattern Request
and Response. In conformità a tale schema asimmetrico, è il client MCP,
su comando dell'host, ad assumere l'iniziativa: esso invia un primo
messaggio di richiesta assegnandogli un identificativo nel campo id (ad
esempio \texttt{1}) ed utilizzando il metodo \texttt{tools/list} per interrogare il
server. Il server restituisce un messaggio di risposta contraddistinto
dal medesimo \texttt{id 1}, nel quale riporta l'elenco dei vari tool disponibili.
Questo include lo strumento \texttt{"get\_weather"}, caratterizzato da un
\texttt{inputSchema}che prescrive l'uso del parametro \texttt{"location"}.

Comprendendo l'obiettivo dell'utente, l'agente AI sceglie di invocare
\texttt{"get\_weather"}. Il client avvia dunque una seconda interazione Request
and Response, inoltrando una nuova richiesta (questa volta con id
impostato a \texttt{2}) avente come metodo \texttt{tools/call}. All'interno di tale
messaggio, il client indica il tool da invocare e i parametri con cui
chiamarlo specificando \texttt{"get\_weather"} nel campo \texttt{name} e passando
l'oggetto \texttt{\{"location": "Torino"\}} nel campo \texttt{arguments},
rispettivamente. Il server esegue l'operazione richiesta e chiude
l'interazione inviando un messaggio di risposta con \texttt{id 2}. Tale messaggio
restituisce l'esito dell'esecuzione all'interno del campo \texttt{content} (ad
esempio, il testo \texttt{"25°C, soleggiato"}) e presenta l'attributo \texttt{isError}
impostato a \texttt{false}, confermando all'agente il successo dell'operazione e
permettendogli di finalizzare la risposta all'utente.

La \zcref{fig:esempio_tool} illustra visivamente l'esempio sopra descritto.

\begin{figure}[htbp]
    \centering
    \includegraphics[width=\textwidth]{diagrams/esempio_tool.pdf}
    \caption{Descrizione dell'interazione tra client e server MCP per l'uso del tool \texttt{get\_weather}.}
    \label{fig:esempio_tool}
\end{figure}

\subsubsection{Risorse}\label{risorse}

Le risorse rappresentano contenuti informativi a cui un agente AI può
accedere, senza modificarle, per acquisire contesto aggiuntivo. Esse
sono identificate in modo univoco tramite URI (Uniform Resource
Identifier) e sono gestite dall'applicazione host. È infatti l'agente AI
a decidere come e quando integrare queste informazioni nel contesto
dell'interazione, anche eventualmente permettendo all'utente di
stabilirlo in modo autonomo.

MCP consente ai server di referenziare le risorse tramite degli URI
fissi che puntano a risorse singole (e.g.~\path{calendar://events/2026}), o
tramite degli URI parametrici, detti resource templates, che possono
essere adattati per richiedere risorse specifiche
(e.g.~\path{weather://forecasts/{location}/{datetime}}). Inoltre, i server
possono supportare l'iscrizione alle risorse, permettendo ai client di
ricevere notifiche in tempo reale se il loro contenuto viene modificato.

\paragraph{Esempio}\label{esempio-risorse}

Per illustrare il funzionamento della primitiva delle risorse, si
considera un agente AI a cui l'utente chiede di riassumere un file di
documentazione locale associato ad un progetto software. L'interazione
per l'accesso a questi dati si articola basandosi nuovamente sul message
pattern Request and Response. Inizialmente, il client inoltra una
richiesta utilizzando il metodo \texttt{resources/list} per scoprire i contenuti
esposti dal server in quel momento. Il server risponde al client
fornendo l'elenco delle risorse disponibili, specificando per ciascuna
di esse un identificatore univoco URI (ad esempio, \path{file:///project/README.md}) e 
il relativo tipo MIME (ad esempio, \texttt{"text/markdown"}).

Una volta che l'utente ha selezionato nell'agente AI la risorsa a cui è
riferita la sua richiesta, il client avvia una seconda interazione
Request and Response inviando una richiesta basata sul metodo
\texttt{resources/read}, all'interno della quale valorizza il parametro \texttt{uri} con
\path{file:///project/README.md}. Il server elabora la chiamata e
restituisce un messaggio di risposta che include il testo effettivo del
documento all'interno di un array denominato \texttt{contents}. A questo punto,
l'agente ha accesso al contenuto e può generare il riassunto per
l'utente.

In aggiunta, essendo i dati delle risorse potenzialmente mutevoli nel
tempo, questa primitiva può esemplificare anche l'utilizzo del message
pattern Subscribe and Notify. In base a tale schema di interazione,
l'agente può inviare tramite il client una richiesta con il metodo
\texttt{subscriptions/listen} indicando l'URI del file README a cui intende
iscriversi. Il server, in risposta, tiene aperto un flusso di
comunicazione continuo su cui emette in modo proattivo una notifica di
tipo \texttt{notifications/resources/updated} ogniqualvolta il file originale
subisce una modifica. Alla ricezione di questo avviso unidirezionale da
parte del server, il client può effettuare in autonomia una nuova
richiesta di lettura, garantendo all'agente un contesto informativo
costantemente sincronizzato con lo stato della risorsa.

La \zcref{fig:esempio_risorse} illustra visivamente l'esempio sopra descritto.

\begin{figure}[htbp]
    \centering
    \includegraphics[width=\textwidth]{diagrams/esempio_risorse.pdf}
    \caption{Descrizione dell'interazione tra client e server MCP per l'uso della risorsa \texttt{file:///project/README.md}.}
    \label{fig:esempio_risorse}
\end{figure}

\subsubsection{Prompt}\label{prompt}

I prompt consistono in template di istruzioni progettati per guidare ed
ottimizzare le interazioni con il modello AI in contesti specifici, ad
esempio per illustrare al LLM con quali tool e/o risorse deve interagire
per svolgere un determinato compito. Più precisamente, un prompt è
articolato in uno o più messaggi per il modello, che possono includere
testo, immagini, audio o collegamenti a risorse specifiche.

MCP supporta inoltre la definizione di prompt parametrici, i quali
possono essere adattati dinamicamente dall'utente durante l'esecuzione
tramite l'inserimento di argomenti. Un possibile esempio di prompt
parametrico è un prompt che richiede al modello AI di revisionare
secondo determinati criteri una porzione di codice che l'utente fornisce
come argomento al momento dell'invocazione.

Infine, si nota che i prompt, a differenza delle risorse, sono gestiti
dagli utenti dell'applicazione AI, che decidono quando invocarli in base
alle loro esigenze.

\paragraph{Esempio}\label{esempio-prompt}

Per illustrare l'utilizzo della primitiva dei prompt, si considera uno
scenario in cui l'utente intende far revisionare una porzione di codice
Python all'agente AI senza occuparsi di formulare la richiesta per il
modello, ma avvalendosi di un template predefinito. L'interazione tra i
componenti si articola seguendo ancora una volta il message pattern
Request and Response. Inizialmente, il client MCP invia una richiesta
utilizzando il metodo \texttt{prompts/list} per ottenere l'insieme dei template
messi a disposizione dal server. Il server restituisce un messaggio di
risposta recante l'elenco dei prompt disponibili, tra cui figura, ad
esempio, un template identificato dal nome \texttt{"code\_review"} che prevede
la fornitura di un parametro ``code'' all'interno dei suoi argomenti.

A differenza dei tool, i prompt sono controllati dall'utente, che li
seleziona in modo esplicito tramite l'interfaccia dell'agente AI. Una
volta che l'utente seleziona il prompt per la revisione e fornisce il
frammento di codice, il client avvia una seconda interazione Request and
Response, inoltrando una richiesta basata sul metodo \texttt{prompts/get}.
All'interno di questo messaggio, il client richiede la risoluzione del
template specificando \texttt{"code\_review"} nel campo \texttt{name} e inserendo il
codice Python all'interno dell'oggetto \texttt{arguments}. Il server processa la
richiesta popolando dinamicamente il template e chiude l'interazione
inviando un messaggio di risposta che include un attributo \texttt{messages}.
Tale campo fornisce un array di messaggi già formattati per il modello
AI, in cui ogni elemento è caratterizzato da un campo \texttt{role} (ad esempio,
\texttt{"user"}) per indicare il mittente e da un campo \texttt{content} contenente le
istruzioni di revisione combinate con il codice Python. L'agente AI
acquisisce così questo contesto pre-strutturato, integrandolo nella
conversazione per soddisfare la richiesta dell'utente.

La \zcref{fig:esempio_prompt} illustra visivamente l'esempio sopra descritto.

\begin{figure}[htbp]
    \centering
    \includegraphics[width=\textwidth]{diagrams/esempio_prompt.pdf}
    \caption{Descrizione dell'interazione tra client e server MCP per l'uso del prompt \texttt{code\_review}.}
    \label{fig:esempio_prompt}
\end{figure}

\subsection{Client}\label{client}

I client MCP operano come intermediari tra l'applicazione host e il
server MCP. Ogni istanza client è creata dall'host per gestire una
connessione con uno specifico server, e si occupa principalmente di
trasferire i messaggi tra host e server in entrambe le direzioni,
allegando ad ogni richiesta i metadati necessari al server per
processarla.

Nell'ultima specifica di MCP, l'unica primitiva che i client possono
esporre è l'elicitation.

\subsubsection{Elicitation}\label{elicitation}

L'elicitation fornisce un meccanismo attraverso il quale i server
possono richiedere informazioni aggiuntive all'utente tramite il client
durante l'elaborazione di una richiesta.

L'approccio adottato dalla primitiva assicura che l'acquisizione delle
informazioni di cui necessita il server avvenga sempre sotto il
controllo diretto dell'applicazione host. In questo modo vengono
garantite trasparenza e sicurezza per l'utente, consentendo inoltre a
quest'ultimo di rifiutare la richiesta del server, se lo desidera.

Dalla specifica 2026-07-28 di MCP, il flusso dell’elicitation segue
il message pattern \hyperref[i-pattern-di-messaggi]{Multi Round-Trip Requests}. 
Il server, avendo bisogno di ulteriori informazioni per soddisfare una 
richiesta del client, replica ad essa con un risultato di tipo \texttt{input\_required}.
Il client che riceve dal server questo tipo di risposta gli inoltra
nuovamente la richiesta a cui è riferita, includendo un campo aggiuntivo
\texttt{inputResponses}. All’interno di tale campo il client specifica se ignora,
respinge o accetta la richiesta di informazioni, e nell’ultimo caso, 
eventualmente, anche un qualche contenuto.


L'elicitation supporta due diverse modalità: form e URL mode.

La \textbf{form mode} consente al server di raccogliere informazioni,
fornendo al client un JSON schema che specifica la struttura dei dati
richiesti all'utente. In questa modalità i dati inseriti dall'utente
sono visibili al client.

La \textbf{URL mode}, invece, permette al server di indirizzare l'utente
verso un URL esterno, dove esso interagisce con un sito web senza che i
dati scambiati vengano intercettati dal client.

Per ragioni di sicurezza, il protocollo impone che i server non
utilizzino la modalità form per richiedere informazioni sensibili, come
password, chiavi API o credenziali di pagamento. Per tali interazioni è
obbligatorio l'uso della modalità URL, che garantisce che dati critici
non transitino attraverso il client MCP.

\paragraph{Esempio}\label{esempio-elicitation}

Al fine di illustrare il funzionamento dell'elicitation, si considera
uno scenario in cui l'agente AI tenta di invocare un tool per
interrogare un sistema esterno, come ad esempio GitHub, ma il server
necessita di conoscere il nome utente per poter procedere. L'interazione
necessaria per acquisire questa informazione esemplifica anche il
message pattern Multi Round-Trip Requests.

Inizialmente, il client MCP inoltra una normale richiesta (assegnandole,
ad esempio, \texttt{1} come identificativo) utilizzando il metodo \texttt{tools/call} per
avviare lo strumento. Il server, constatando l'assenza del nome utente
necessario per operare, sospende l'esecuzione e restituisce una risposta
contenente un risultato di tipo \texttt{input\_required}. All'interno di tale
messaggio, il server inserisce un oggetto \texttt{inputRequests} che incapsula
una richiesta \texttt{elicitation/create}. Per raccogliere il dato, il server si
avvale della form mode: imposta dunque il parametro mode a \texttt{"form"} e
fornisce un \texttt{requestedSchema} in formato JSON Schema, che prescrive al
client di richiedere una stringa di testo corrispondente al nome utente.

L'applicazione host, interpretando la richiesta, presenta un modulo
all'utente, il quale inserisce il proprio nome (ad esempio, \texttt{"Paolo58"})
e autorizza la condivisione. A questo punto, in aderenza alla logica del
pattern Multi Round-Trip Requests, il client avvia il secondo ciclo
dell'interazione (round-trip) reiterando la richiesta \texttt{tools/call}
originale, alla quale però assegna un nuovo identificativo univoco (ad
esempio \texttt{2}). In questa nuova chiamata, il client integra il campo
aggiuntivo \texttt{inputResponses}, all'interno del quale specifica l'avvenuta
autorizzazione impostando \texttt{action} ad \texttt{"accept"} e fornisce il dato
acquisito tramite l'oggetto \texttt{\{"name": "Paolo58"\}} nel campo \texttt{content}.
Ricevute le informazioni mancanti, il server riprende l'elaborazione
della chiamata e conclude l'interazione inviando un messaggio di
risposta finale, associato all'\texttt{id 2}, contenente l'esito dell'esecuzione
del tool.

La \zcref{fig:esempio_elicitation} illustra visivamente l'esempio sopra descritto.

\begin{figure}[htbp]
    \centering
    \includegraphics[width=\textwidth]{diagrams/esempio_elicitation.pdf}
    \caption{Descrizione dell'interazione tra client e server MCP per l'uso di un tool che richiede informazioni aggiuntive tramite elicitation.}
    \label{fig:esempio_elicitation}
\end{figure}

\section{Il Blindingly Simple Protocol Language}\label{bspl}

Nei sistemi informativi distribuiti composti da componenti software eterogenei e autonomi, 
spesso sviluppati da organizzazioni differenti e operanti su piattaforme non omogenee, la 
cooperazione richiede la definizione di regole formali di comunicazione condivise. 
L'\emph{Interaction-Oriented Programming} (IOP) affronta questo scenario spostando il fulcro 
della progettazione dalla struttura interna dei singoli componenti alla specifica delle 
loro interazioni pubbliche \autocite{singh2011bspl}.

All'interno di questo paradigma, un \textbf{protocollo} definisce lo schema formale delle 
comunicazioni consentite tra determinati \textbf{ruoli}. In tal modo, ciascun partecipante 
mantiene la propria autonomia implementativa e decisionale, con l'unico vincolo di conformarsi 
alla grammatica dei messaggi stabilita dal protocollo durante l'interazione.

Tradizionalmente, la modellazione dei protocolli di interazione nei sistemi multi-agente è 
stata affrontata impiegando formalismi di tipo procedurale, quali i \emph{Message Sequence 
Charts} (MSC), AUML, BPMN e WS-CDL. Questi linguaggi descrivono l'interazione attraverso diagrammi 
di flusso, prescrivendo in modo esplicito sequenze rigide di messaggi, 
diramazioni condizionali e costrutti di sincronizzazione imperativi.

Tuttavia, l'imposizione di un flusso di controllo procedurale presenta criticità rilevanti 
quando applicata a contesti distribuiti ed asincroni. Da un lato, tali modelli tendono a 
risultare sovra-vincolati, poiché stabiliscono un ordine temporale 
rigido anche tra comunicazioni che non presentano alcuna reale dipendenza logica; ciò limita 
il parallelismo e costringe gli sviluppatori a realizzare adattatori complessi per i singoli 
agenti, aumentando l'accoppiamento del sistema. Dall'altro lato, in reti prive di un'autorità 
centrale o di una memoria condivisa, i flussi di controllo procedurali faticano a gestire 
l'asincronia, esponendo il sistema ad anomalie di coordinamento quali la \emph{nonlocal choice} 
(in cui decisioni concorrenti ed incompatibili, prese in modo indipendente da agenti distinti, 
generano stati di inconsistenza o di stallo) e la \emph{blindness} (in cui un ruolo non ha 
visibilità locale sufficiente per verificare se un vincolo di precedenza globale sia stato 
effettivamente rispettato) \autocite{singh2011lost}.

Per superare la rigidità dei modelli procedurali, la ricerca nel settore dei sistemi 
multi-agente ha proposto approcci dichiarativi incentrati sul significato di business delle 
interazioni, ad esempio modellando i messaggi in termini di impegni sociali 
tra le parti. Ciononostante, per essere eseguiti in pratica, questi approcci semantici 
necessitavano comunque di regole operazionali precise per disciplinare l'invio effettivo dei 
messaggi, finendo spesso per riutilizzare vincoli procedurali sull'ordine delle comunicazioni e 
reintroducendo i medesimi problemi di accoppiamento.

Proprio con l'obiettivo di definire regole di esecuzione concrete ma puramente dichiarative, nel 2011 
Munindar P. Singh ha introdotto il \textbf{Blindingly Simple Protocol Language} 
(\textbf{BSPL}) \autocite{singh2011bspl}.

In BSPL, il protocollo viene trattato come una \emph{first-class abstraction}: questo significa
che l'interazione è modellata come un modulo autonomo, incapsulato e dotato di una propria 
interfaccia pubblica (definita da ruoli e parametri), il quale può essere specificato, verificato 
formalmente e riutilizzato per composizione con altri protocolli senza dover dipendere dal codice 
interno degli agenti o da motori di orchestrazione centralizzati.

L'elemento distintivo di BSPL risiede nel suo approccio guidato dall'informazione: il linguaggio 
non contiene costrutti per imporre sequenze, bivi o cicli tra i messaggi. L'ordine in cui le 
comunicazioni possono avvenire non viene fissato a priori nella struttura di controllo, ma scaturisce 
proprio dalle informazioni che i messaggi trasmettono, più precisamente dalle dipendenze causali 
tra queste informazioni.

\subsection{Principi guida}\label{bspl:principi}

La formulazione di BSPL poggia su quattro principi fondamentali che ne guidano sia la 
struttura che la semantica operazionale \autocite{singh2011bspl}:

Innanzitutto, il principio dell'\textbf{orientamento all'informazione} stabilisce che la finalità 
principale per cui gli agenti interagiscono risieda nello scambio e nella generazione di dati.
Di conseguenza, un protocollo BSPL specifica, accanto ai ruoli partecipanti, l'insieme dei 
\textbf{parametri} che ne costituiscono il contenuto informativo. Il progresso dell'interazione 
coincide operativamente con la progressiva assegnazione di valori (\emph{binding}) a tali parametri.

In secondo luogo, la \textbf{causalità esplicita} stabilisce che 
l'ordinamento temporale tra le comunicazioni dipenda unicamente dalla disponibilità delle 
informazioni. In BSPL, un agente è abilitato a inviare un messaggio solo nel 
momento in cui possiede localmente tutti i parametri richiesti in ingresso da quello specifico 
schema di messaggio. Pertanto, se un messaggio necessita di un dato prodotto esclusivamente da 
un altro messaggio, il primo potrà essere emesso solo dopo che il secondo ha avuto luogo e il 
dato è stato comunicato. Al contrario, se due messaggi non presentano dipendenze informative 
reciproche, gli agenti possono inviarli in qualsiasi ordine o simultaneamente, massimizzando 
la concorrenza ed evitando vincoli di sequenziamento artificiosi.

Inoltre, il modello adotta il principio dell'\textbf{assenza di stato globale}, 
non richiedendo alcun database condiviso o sincronizzazione centrale 
tra i nodi. L'intero stato pubblico dell'interazione (definito \emph{social state}) è 
tracciabile esclusivamente nei messaggi scambiati. Ciascun agente determina la 
legittimità delle proprie azioni basandosi unicamente sul proprio stato locale (\emph{local 
state}), ovvero sull'insieme dei dati appresi attraverso i messaggi che ha personalmente 
inviato o ricevuto.

Infine, BSPL garantisce la netta \textbf{separazione tra struttura operazionale e significato}.
Il protocollo si occupa della correttezza operazionale delle comunicazioni, 
stabilendo chi produce e consuma ciascun parametro e quando un messaggio sia ammissibile.
Si astiene, invece, dal vincolare la logica decisionale privata con cui ciascun agente elabora
i propri dati internamente o il valore semantico delle transazioni. 
Tale disaccoppiamento consente di aggiungere formalismi semantici di livello superiore, 
come gli impegni contrattuali, su una base operazionale distribuita.

\subsection{Modello informativo}\label{bspl:modello}

BSPL formalizza l'interazione tra agenti adottando un approccio che prende spunto 
dal modello relazionale di Codd \autocite{singh2011bspl}. In questo formalismo, 
un protocollo viene trattato come uno schema di relazione, composto dai parametri definiti all'interno
del protocollo stesso, mentre ciascuna sua singola esecuzione (\emph{enactment}) corrisponde a una tupla di 
valori assegnati a tali parametri.

Più precisamente, una specifica di protocollo in BSPL è costituita da quattro elementi fondamentali:
\begin{itemize}
    \item L'insieme $\mathcal{R}$ dei \textbf{ruoli}, che identificano le entità logiche 
    coinvolte nell'interazione in qualità di mittenti o destinatari dei messaggi.
    \item L'insieme $\mathcal{P}$ dei \textbf{parametri}, che rappresentano le variabili 
    informative scambiate durante l'esecuzione.
    \item Un sottoinsieme $\vec{k} \subseteq \mathcal{P}$ di parametri contrassegnati come \textbf{chiavi}, deputati a distinguere univocamente le diverse 
    istanze di esecuzione.
    \item Un insieme $\mathcal{M}$ di riferimenti a \textbf{schemi di messaggio} o a 
    sottoprotocolli che articolano l'interazione complessiva.
\end{itemize}

\subsubsection{Assiomi fondativi}\label{bspl:assiomi}

La correttezza delle esecuzioni distribuite in BSPL si fonda su tre 
assiomi riguardanti la gestione dei parametri e delle loro istanze:

Innanzitutto, l'\textbf{unicità} impone che ogni esecuzione di un protocollo sia 
identificata in modo univoco dai valori assegnati ai parametri che ne costituiscono la 
chiave. A livello operativo, questo meccanismo garantisce che messaggi 
recapitati in momenti differenti o attraverso canali asincroni distinti possano essere 
associati con certezza alla medesima sessione di interazione. Si evita così qualsiasi 
interferenza tra esecuzioni concorrenti senza la necessità di stabilire canali dedicati a 
livello di trasporto o di utilizzare identificatori di sessione ad-hoc.

In secondo luogo, l'\textbf{integrità} stabilisce che un'istanza di esecuzione sia 
considerata completa solo se a tutti i parametri pubblici previsti dall'interfaccia 
del protocollo è stato associato un valore concreto. In modo analogo ai vincoli 
\texttt{NOT NULL} delle basi di dati relazionali, la presenza di parametri non ancora 
vincolati indica che l'interazione è parziale e necessita di ulteriori scambi di messaggi 
per giungere a compimento.

Infine, l'\textbf{immutabilità} prescrive che, all'interno di una specifica esecuzione,
una volta che a un parametro è stato assegnato un valore, 
tale assegnazione non possa più essere modificata o sovrascritta. Questo vincolo ha 
ricadute ingegneristiche fondamentali per i sistemi distribuiti: rende l'elaborazione dei 
messaggi intrinsecamente idempotente ed immune da \emph{race condition} sui dati. Qualora 
un messaggio venga duplicato a causa di ritrasmissioni su canali inaffidabili o venga 
recapitato fuori ordine su canali non-FIFO, l'agente ricevente è in grado di rilevare la 
ridondanza e preservare la consistenza del proprio stato locale senza incorrere in conflitti 
o ambiguità.

\paragraph{Esempio}\label{bspl:esempio_modello_relazionale}
Per illustrare come BSPL applichi il modello relazionale all'esecuzione distribuita, 
si può considerare un protocollo di acquisto (\zcref{lst:bspl_purchase}) caratterizzato dai parametri \texttt{ID}
(che rappresenta la chiave), \texttt{item}, \texttt{price} e \texttt{receipt}.
Ciascun agente mantiene localmente una tabella in cui ogni riga corrisponde a una specifica istanza
di interazione, identificata in modo univoco dal valore del parametro chiave \texttt{ID}.

Durante lo scambio di messaggi, la tupla associata a una determinata sessione non subisce sovrascritture 
distruttive, ma evolve in modo puramente monotono crescente attraverso il progressivo assegnamento dei parametri.
La \zcref{tab:bspl_relational_example} illustra lo stato locale di un agente con ruolo \textsc{Buyer} in un
dato momento.

\begin{table}[htbp]
\centering
\caption{Vista relazionale di un esempio di stato locale per il ruolo \textup{\textsc{Buyer}}.}
\label{tab:bspl_relational_example}
\begin{tabularx}{\textwidth}{l l l l >{\raggedright\arraybackslash}X}
\toprule
\textbf{Istanza (\texttt{ID})} & \textbf{\texttt{item}} & \textbf{\texttt{price}} & \textbf{\texttt{receipt}} & \textbf{Stato dell'enactment} \\
\midrule
\texttt{\#100} & \texttt{"Mouse"}    & 25\,€     & \texttt{"REC-112"} & \emph{Completo} (tupla integra) \\
\texttt{\#101} & \texttt{"Libro"}    & 20\,€     & $\bot$             & \emph{Parziale} (in attesa di ricevuta) \\
\texttt{\#102} & \texttt{"Tastiera"} & $\bot$    & $\bot$             & \emph{Parziale} (in attesa di offerta) \\
\bottomrule
\end{tabularx}
\end{table}

Considerando i tre assiomi su cui si regge BSPL in relazione alla \zcref{tab:bspl_relational_example} si può
innanzitutto osservare che ciascuna riga rappresenta un'esecuzione indipendente del protocollo,
distinta univocamente dalla propria chiave (\texttt{\#100}, \texttt{\#101}, \texttt{\#102}). 
Inoltre, vale che quando un dato, come il prezzo di 20\,€ per l'istanza \texttt{\#101}, viene appreso 
e inserito nella tabella, esso non può più essere modificato. L'evoluzione della sessione consiste
esclusivamente nel sostituire il valore indefinito ($\bot$) dei parametri non ancora vincolati. 
Infine, si nota che l'istanza \texttt{\#100} rappresenta un'esecuzione conclusa con successo poiché tutti 
i parametri pubblici previsti dal protocollo risultano vincolati, mentre le istanze \texttt{\#101}
e \texttt{\#102} descrivono sessioni ancora incomplete che necessitano di ulteriori scambi di messaggi 
per giungere a compimento.

\subsection{Sintassi}\label{bspl:sintassi}

La sintassi di BSPL è concepita per essere minimale e priva di costrutti espliciti di controllo del flusso,
come bivi o loop. Essa si articola principalmente nella caratterizzazione dei parametri, nella definizione 
di schemi di messaggio e nella struttura di dichiarazione modulare del protocollo, suddivisa in intestazione
e corpo.

\subsubsection{Gli adornamenti dei parametri}\label{bspl:adornamenti}

In assenza di costrutti espliciti per il controllo del flusso, BSPL governa la causalità e 
la validità delle comunicazioni mediante gli \textbf{adornamenti} associati ai parametri. 
Gli adornamenti compaiono sia all'interno dei singoli schemi di messaggio, sia nell'intestazione 
delle dichiarazioni di protocollo.

All'interno degli \textbf{schemi di messaggio}, il linguaggio definisce tre tipologie di adornamento.

Innanzitutto, si ha l'adornamento \textbf{\adorn{in}}, il quale indica che il parametro 
costituisce un prerequisito conoscitivo per l'invio del messaggio. Per poter emettere un
messaggio contenente un parametro \adorn{in}, il ruolo mittente
deve già conoscere il valore di tale parametro avendolo appreso da una ricezione 
precedente o avendolo generato in un messaggio emesso in precedenza.

L'adornamento \textbf{\adorn{out}} indica invece che il parametro 
rappresenta un'informazione generata per la prima volta da quel 
messaggio. L'invio di un messaggio con un parametro \adorn{out} possiede 
forza dichiarativa: il mittente non si limita a trasmettere un valore, ma lo fissa in 
modo definitivo per l'istanza corrente del protocollo. Da quel momento, il mittente, e 
successivamente il destinatario, conoscono il dato, e nessun ruolo potrà emettere altri 
messaggi che dichiarano lo stesso parametro con \adorn{out} per la 
medesima chiave.

Infine, l'adornamento \textbf{\adorn{nil}} specifica che il parametro
deve essere non vincolato nello stato di conoscenza locale del mittente al momento 
dell'emissione del messaggio.
A differenza della semplice omissione di un parametro dallo schema del messaggio, per cui 
la sua presenza o assenza nella conoscenza locale è ininfluente ai fini dell'invio, 
l'uso di \adorn{nil} impedisce l'emissione qualora il mittente abbia già appreso o generato 
tale informazione. In questo modo, il protocollo può stabilire una mutua esclusione 
tra schemi di messaggio alternativi o subordinare l'invio alla mancata conoscenza di un dato.

Nell'\textbf{intestazione di un protocollo}, gli adornamenti definiscono invece il contratto 
d'interfaccia verso l'esterno. Un parametro dichiarato con \adorn{in} rappresenta 
un dato che deve essere fornito dall'esterno (da un altro protocollo o dall'ambiente) affinché 
l'interazione possa svolgersi.
Invece, un parametro dichiarato con \adorn{out} rappresenta un dato che il protocollo
si impegna a generare e rendere disponibile all'esterno al termine dell'esecuzione.

L'intestazione impone due vincoli di buona formazione sui messaggi definiti nel corpo del protocollo:
un parametro dichiarato \adorn{in} nell'intestazione può comparire esclusivamente 
come \adorn{in} nei messaggi interni, poiché non può essere generato all'interno 
del protocollo stesso; al contrario, un parametro dichiarato \adorn{out} nell'intestazione 
deve essere generato con \adorn{out} da almeno uno schema di messaggio interno.

Questa distinzione consente di identificare i \textbf{protocolli top-level}, ovvero protocolli 
autonomi e completi, pronti per essere eseguiti direttamente da un sistema multi-agente senza 
dipendere da informazioni esterne. Un protocollo top-level si caratterizza per avere tutti i 
suoi parametri pubblici adornati con \adorn{out}, inclusa la chiave dell'interazione 
generata dal ruolo che dà inizio all'interazione. Al contrario, un protocollo che presenta 
parametri \adorn{in} nella propria intestazione costituisce un frammento non 
autosufficiente, il quale può essere eseguito solo all'interno di una composizione che ne 
generi i prerequisiti.

\subsubsection{La visibilità dei parametri}\label{bspl:visibilita}

Oltre all'adornamento, BSPL disciplina la \textbf{visibilità} dei parametri distinguendo tra
parametri pubblici e locali.

I \textbf{parametri pubblici} sono quelli dichiarati esplicitamente nella clausola \texttt{parameter} 
dell'intestazione del protocollo. Essi definiscono l'interfaccia visibile all'esterno e concorrono 
a stabilire la condizione di completamento dell'esecuzione.

I \textbf{parametri locali}, invece, non richiedono alcuna parola chiave per essere dichiarati: sono 
semplicemente parametri che compaiono all'interno dei messaggi nel corpo del protocollo ma vengono 
omessi dalla lista dei parametri nell'intestazione. Nei messaggi, i parametri locali vengono 
adornati con \adorn{in}, \adorn{out} o \adorn{nil} 
esattamente come quelli pubblici.

I parametri locali svolgono due funzioni fondamentali. 
In primo luogo consentono ai ruoli partecipanti 
di scambiare informazioni necessarie allo svolgimento dell'interazione (come un indirizzo di spedizione 
\texttt{address} o un codice di tracciamento) senza esporle nell'interfaccia pubblica del protocollo, 
preservando l'astrazione e il disaccoppiamento verso eventuali protocolli chiamanti. 
In secondo luogo, consentono di modellare percorsi alternativi senza compromettere il completamento dell'istanza.
In un bivio decisionale due messaggi alternativi come \texttt{accept} e \texttt{reject} possono
produrre ciascuno informazioni diverse, come un codice di tracciamento o una motivazione di scarto.
Poiché i parametri locali non appartengono all'interfaccia pubblica, il fatto che alcuni di essi rimangano
non assegnati al termine dell'interazione non impedisce all'istanza del protocollo di raggiungere lo stato di
completamento, il quale richiede unicamente l'assegnazione di tutti i parametri pubblici.

\subsubsection{Il modificatore \texttt{key}}\label{bspl:key}
Accanto agli adornamenti \texttt{in}, \texttt{out} e \texttt{nil}, BSPL introduce il modificatore \texttt{key}, il quale può essere
utilizzato esclusivamente nell'intestazione del protocollo per qualificare uno o più parametri come chiave 
primaria del suo schema relazionale.

Come già accennato nel \zcref{bspl:assiomi}, la chiave di un protocollo è l'insieme dei parametri 
che identificano in modo univoco ciascuna sua istanza di esecuzione: tutti i messaggi che presentano 
la medesima combinazione di valori per i parametri chiave appartengono alla stessa sessione del protocollo.

Si nota, inoltre, che uno schema di messaggio deve obbligatoriamente contenere almeno un parametro chiave del protocollo di cui fa parte,
ma non necessariamente tutti. Quando un protocollo definisce una chiave composta da più parametri, i messaggi che
includono soltanto una parte della chiave fissano informazioni comuni a tutte le istanze del protocollo che condividono quel valore parziale.
Nel momento in cui l'interazione deve produrre o consumare informazioni specifiche di una singola sessione, lo schema
del messaggio deve necessariamente includere anche i restanti parametri chiave discriminanti.

\subsubsection{Schemi di messaggio}\label{bspl:schemi-messaggio}

Definiti i ruoli e la qualificazione dei parametri, l'unità fondamentale di 
comunicazione in BSPL è costituita dallo schema di messaggio.
Uno schema di messaggio rappresenta un'interazione atomica tra due ruoli, un mittente e un 
destinatario, e assume la forma:
\begin{center}
    \bsplinline{Mittente -> Destinatario: NomeMessaggio[adornamento Parametro1, ...]}
\end{center}
in cui ciascun parametro è specificato con il relativo adornamento (\adorn{in}, 
\adorn{out} o \adorn{nil}).

\subsubsection{Dichiarazione di un protocollo}\label{bspl:dichiarazione-protocollo}

La dichiarazione di un protocollo definisce un blocco racchiuso tra parentesi graffe contenente 
il nome del protocollo, i ruoli partecipanti (\texttt{role}), i parametri pubblici (\texttt{parameter}) 
e il corpo del protocollo, costituito da schemi di messaggio o riferimenti ad altri protocolli.

La \zcref{lst:bspl_offer} e la \zcref{lst:bspl_pay} illustrano due protocolli elementari: 
\emph{Offer}, in cui il ruolo \textsc{Buyer} richiede un preventivo e il ruolo \textsc{Seller} 
produce il prezzo, e \emph{Pay}, in cui il ruolo \textsc{Payer} trasferisce un importo al ruolo \textsc{Payee}.

\begin{bspllisting}[caption={Dichiarazione del protocollo Offer in BSPL.}, label={lst:bspl_offer}]
Offer {
    role Buyer, Seller
    parameter in ID key, out item, out price

    Buyer -> Seller: rfq[in ID, out item]
    Seller -> Buyer: quote[in ID, in item, out price]
}
\end{bspllisting}

\begin{bspllisting}[caption={Dichiarazione del protocollo Pay in BSPL.}, label={lst:bspl_pay}]
Pay {
    role Payer, Payee
    parameter in ID key, in amount, out paid

    Payer -> Payee: payM[in ID, in amount, out paid]
}
\end{bspllisting}

\subsubsection{Composizione di protocolli}\label{bspl:composizione}

Il meccanismo di \textbf{composizione} consente di riutilizzare protocolli esistenti all'interno 
di una nuova specifica mediante riferimenti espliciti nel corpo della dichiarazione. 
La sintassi di un riferimento a un sottoprotocollo assume la forma:
\begin{center}
    \bsplinline{NomeSottoprotocollo(Ruolo1, ..., adornamento Parametro1, ...)}
\end{center}

All'interno della composizione i ruoli e i parametri del sottoprotocollo vengono mappati 
rispettivamente sui ruoli e sui parametri (pubblici o locali) del protocollo chiamante.
Poiché in BSPL l'identità di un'istanza di un protocollo è determinata unicamente dai valori delle sue chiavi,
la grammatica impone che il sottoprotocollo condivida almeno un parametro chiave con il protocollo
chiamante. Tale vincolo funge da meccanismo di correlazione, garantendo che i messaggi scambiati
nel sottoprotocollo siano univocamente associati alla medesima sessione di esecuzione del protocollo principale.

La \zcref{lst:bspl_purchase} mostra un esempio concreto di composizione: il protocollo top-level 
\emph{Purchase} integra i sottoprotocolli \emph{Offer} e \emph{Pay}, completando l'acquisto con
l'invio di una ricevuta per l'acquisto da parte del \textsc{Seller} tramite il messaggio \texttt{receiptM}.

\begin{bspllisting}[caption={Dichiarazione del protocollo composito Purchase in BSPL.}, label={lst:bspl_purchase}]
Purchase {
    role Buyer, Seller
    parameter out ID key, out item, out price, out receipt

    Buyer -> Seller: init[out ID]
    Offer(Buyer, Seller, in ID, out item, out price)
    Pay(Buyer, Seller, in ID, in price, out paid)
    Seller -> Buyer: receiptM [in ID, in price, in paid, out receipt]
}
\end{bspllisting}

Nella \zcref{lst:bspl_purchase}, il protocollo \emph{Purchase} associa il parametro \texttt{price} 
generato come \adorn{out} da \emph{Offer} al parametro \texttt{amount} 
richiesto come \adorn{in} da \emph{Pay}. 
In questo modo, la corretta sequenza operativa tra l'offerta e il pagamento scaturisce 
automaticamente e unicamente dalla disponibilità del dato \texttt{price}: \emph{Pay} non potrà 
essere eseguito prima che \emph{Offer} abbia prodotto tale valore. La composizione rispetta 
l'incapsulamento, poiché \emph{Purchase} interagisce con i sottoprotocolli esclusivamente attraverso 
le loro interfacce pubbliche.

\subsection{Semantica operazionale}\label{bspl:semantica}

Nei sistemi distribuiti asincroni non esiste una memoria globale condivisa né un'autorità centrale
in grado di coordinare i partecipanti. Di conseguenza, la semantica operazionale di BSPL descrive
l'esecuzione di un protocollo dal punto di vista decentralizzato dei singoli agenti, basandosi
esclusivamente sulle informazioni apprese localmente da ciascun ruolo durante lo scambio di messaggi \autocite{singh2012semantics}.

\subsubsection{Storia locale e stato di conoscenza locale}\label{bspl:stato_locale}

Durante l'esecuzione, un agente che ricopre un determinato ruolo $x \in \mathcal{R}$ non ha visibilità
sulle comunicazioni che avvengono tra gli altri partecipanti. L'agente percepisce l'interazione unicamente
attraverso gli eventi comunicativi a cui prende parte direttamente. L'insieme dei messaggi che il ruolo
$x$ ha personalmente inviato o ricevuto costituisce la sua \textbf{storia locale} $\mathcal{H}_x$.

A partire dalla propria storia locale, l'agente costruisce il proprio \textbf{stato di conoscenza locale}
$\mathcal{K}_x$. Per ciascuna istanza di interazione, identificata dal valore della chiave $\vec{k}$,
lo stato $\mathcal{K}_x$ memorizza i valori assegnati ai parametri che il ruolo ha appreso fino a quel
momento. Se un parametro non è ancora comparso in alcun messaggio inviato o ricevuto da $x$ per la
chiave $\vec{k}$, il suo valore per quell'agente rimane indefinito ($\bot$).

In questo modo, la valutazione della liceità delle azioni non richiede alcuna sincronizzazione globale:
ciascun agente determina in modo del tutto autonomo quali messaggi è abilitato a inviare o può attendersi
di ricevere basandosi unicamente sulle informazioni presenti nel proprio stato di conoscenza locale.

\subsubsection{Regole di viabilità per emissione e ricezione dei messaggi}\label{bspl:viabilita}

Il concetto cardine che disciplina l'invio e la ricezione dei messaggi in BSPL è la 
\textbf{viabilità} (\emph{viability}), la quale stabilisce se un evento comunicativo sia 
formalmente consentito dallo stato di conoscenza locale dell'agente.

L'\textbf{emissione} di un messaggio $m$ da parte di un ruolo mittente $x$ verso un destinatario $y$, 
relativamente a una sessione identificata dalla chiave $\vec{k}$, è viabile se e solo se sono 
soddisfatte simultaneamente due condizioni su $\mathcal{K}_x$:
\begin{enumerate}
    \item Il mittente $x$ deve già conoscere il valore 
    di tutti i parametri adornati con \adorn{in} nello schema di $m$ per la 
    chiave $\vec{k}$, ovvero tali parametri devono essere vincolati in $\mathcal{K}_x$.
    \item Il mittente $x$ non deve conoscere alcun valore 
    per i parametri adornati con \adorn{out} o con \adorn{nil} 
    nello schema di $m$ per la chiave $\vec{k}$, ossia tali parametri devono risultare non vincolati in $\mathcal{K}_x$.
\end{enumerate}

La prima condizione assicura il rispetto delle dipendenze causali: un agente non può trasmettere un dato che non possiede. 
La seconda condizione preserva l'assioma di immutabilità: se $x$ conoscesse già un parametro \adorn{out}, 
l'invio del messaggio costituirebbe una riassegnazione incompatibile con l'unicità del valore nella sessione.

Nel momento in cui avviene l'invio del messaggio, $x$ genera i valori per ciascun parametro \adorn{out} 
e li inserisce immediatamente in $\mathcal{K}_x$. Da questo istante, la seconda condizione di viabilità 
impedirà qualsiasi altro invio da parte di $x$ che tenti di riassegnare quegli stessi 
parametri per la chiave $\vec{k}$.

La \textbf{ricezione} di un messaggio $m$ da parte del destinatario $y$ è invece un'operazione \textbf{incondizionata},
ovvero sempre viabile a livello di semantica formale.
All'atto della ricezione, $y$ incorpora i valori contenuti nel messaggio all'interno del proprio stato locale 
$\mathcal{K}_y$, vincolando i parametri che prima risultavano indefiniti. L'arricchimento di $\mathcal{K}_y$ può 
quindi abilitare a sua volta l'emissione di successivi messaggi da parte di $y$, qualora i parametri ricevuti soddisfino 
i requisiti \adorn{in} di altri schemi di messaggio.

\subsubsection{Monotonicità della conoscenza}\label{bspl:monotonicita}

Una conseguenza importante delle regole di viabilità e dell'assioma di immutabilità è la 
\textbf{monotonicità} dell'evoluzione dello stato locale: sia l'invio che la ricezione di un messaggio 
possono unicamente aggiungere nuove informazioni alla conoscenza dell'agente, senza mai sovrascrivere o 
cancellare i dati già appresi.

La crescita monotona garantisce che le decisioni locali degli agenti siano stabili nel tempo. 
Una volta che un agente ha appreso i parametri \adorn{in} necessari ad abilitare un messaggio, 
tale condizione di abilitazione non può essere revocata da eventi asincroni successivi.

Questo modello offre due importanti vantaggi architetturali.
Innanzitutto, gli agenti possono operare in modo del tutto autonomo e disaccoppiato nel tempo. 
Poiché lo stato di conoscenza locale non rischia di essere invalidato da azioni 
concorrenti, un ruolo è autorizzato a emettere un messaggio nel momento esatto 
in cui questo diviene viabile localmente, senza dover implementare protocolli 
di sincronizzazione preventiva, negoziare permessi o acquisire blocchi 
sullo stato degli altri nodi. 
Inoltre, nei sistemi transazionali tradizionali, la concorrenza 
e la gestione delle eccezioni richiedono costosi protocolli di coordinamento 
globale, come il \emph{Two-Phase Commit}, o complesse logiche di ripristino per annullare 
mutazioni di stato inconsistenti. In BSPL, l'esecuzione modella l'interazione come la progressiva 
popolazione di una tupla relazionale. Anche la gestione di esiti negativi o rami di 
scarto (come un messaggio di \emph{reject}) non cancella i dati precedentemente 
stabiliti, ma progredisce vincolando parametri previsti dal protocollo. 
Evolvendo l'interazione unicamente per accumulo di informazione, si elimina 
alla radice la necessità di ritrattare comunicazioni o ripristinare stati pregressi.

\subsection{Architettura dell'agente}\label{bspl:architettura_agente}

La semantica operazionale di BSPL si traduce a livello implementativo in una netta separazione 
delle responsabilità all'interno di ciascun nodo partecipante. Nel modello architetturale Local State Transfer
(LoST) \autocite{singh2011lost}, un agente non è strutturato come un blocco monolitico, 
ma si articola in due componenti disaccoppiati: l'\textbf{adattatore di protocollo} e la \textbf{logica di business}.

\subsubsection{Adattatore}\label{bspl:adattatore}

L'adattatore di protocollo agisce come uno strato di \emph{middleware} interposto tra 
l'infrastruttura di comunicazione e l'applicazione dell'agente. 
Questo componente opera in modo generico e indipendente dallo specifico dominio applicativo,
avendo come responsabilità principale quella di gestire l'invio e la ricezione dei messaggi
sul canale di trasporto, aggiornando lo stato di conoscenza locale. 
L'adattatore si occupa, inoltre, di applicare le regole di viabilità, impedendo l'emissione 
di messaggi privi dei necessari prerequisiti \adorn{in} o in conflitto con
parametri \adorn{out} già vincolati. 
Infine, esso espone un'interfaccia verso l'interno dell'agente, notificando l'acquisizione di
nuove informazioni o l'abilitazione di nuove comunicazioni.

\subsubsection{Logica di business}\label{bspl:logica}

La logica di business, al contrario, costituisce il modulo privato dell'agente. 
Essa racchiude i processi decisionali e gli algoritmi di calcolo, 
avendo come unica responsabilità la generazione dei valori da assegnare ai parametri \adorn{out}. 
La logica interna non gestisce la conformità al protocollo né le dinamiche di trasporto: quando 
intende compiere un'azione comunicativa, delega l'emissione all'adattatore fornendogli i dati elaborati.

Questo disaccoppiamento garantisce che l'adattatore sia un modulo interamente riutilizzabile su 
qualsiasi specifica BSPL. Inoltre, protegge il sistema distribuito da malfunzionamenti interni: 
se la logica di business tenta un'azione non consentita, l'adattatore la blocca localmente, 
evitando che violazioni del protocollo si propaghino sulla rete.


\subsubsection{Struttura dei dati}\label{bspl:struttura-dati}

Per gestire lo stato di conoscenza locale, l'adattatore non mantiene una tabella globale 
del protocollo. Al contrario, memorizza una \textbf{relazione locale} $R(m)$ 
per ciascuno schema di messaggio $m$ in cui il ruolo partecipa come mittente o destinatario \autocite{singh2011lost}.
In questo modo, lo stato di conoscenza corrisponde all'insieme delle tuple
memorizzate nelle relazioni locali gestite dall'adattatore.

È opportuno evidenziare la distinzione tra la vista logica del protocollo e la sua materializzazione fisica all'interno dell'adattatore. Se
a livello concettuale lo stato dell'interazione viene descritto come un'unica tabella che raccoglie tutti i parametri per ciascuna
sessione, a livello implementativo l'adattatore non memorizza direttamente tale relazione globale. Una tabella unificata, infatti,
costringerebbe l'adattatore a eseguire continui aggiornamenti distruttivi per rimpiazzare i valori indefiniti a ogni nuova
informazione appresa, oltre a risultare inadatta a gestire i rami alternativi o i messaggi a cui il ruolo non prende parte. Al contrario,
materializzare una relazione locale dedicata per ciascuno schema di messaggio consente di gestire lo stato esclusivamente tramite puri
inserimenti di tuple immutabili, rispecchiando la natura incrementale della conoscenza. In questa architettura, la tabella complessiva del 
protocollo rappresenta unicamente una vista logica derivata, ottenibile combinando le singole relazioni locali attraverso le relative chiavi
di sessione.

\paragraph{Esempio}

Per comprendere come la vista logica sullo stato dell'interazione introdotta nella \zcref{tab:bspl_relational_example} (riferita all'esempio nel \zcref{bspl:esempio_modello_relazionale})
venga materializzata a livello fisico, si consideri il medesimo protocollo di acquisto (\zcref{lst:bspl_purchase}) dal punto di vista dell'adattatore del ruolo \textsc{Buyer}.
L'interazione per \textsc{Buyer} si articola negli schemi di messaggio seguenti: l'invio dell'inizializzazione \bsplinline{init[out ID]}, l'invio della richiesta \bsplinline{rfq[in ID, out item]},
la ricezione del preventivo \bsplinline{quote[in ID, in item, out price]}, l'invio del pagamento \bsplinline{payM[in ID, in price, out paid]} e la ricezione della ricevuta finale 
\bsplinline{receiptM[in ID, in price, in paid, out receipt]}.

L'adattatore di \textsc{Buyer} non memorizza una tabella sparsa con valori indefiniti, ma materializza cinque diverse relazioni locali,
una per ogni schema di messaggio in cui è coinvolto, illustrate nella \zcref{tab:bspl_physical_example}.

\begin{table}[htbp]
\centering
\small
\caption{Materializzazione fisica delle relazioni locali nell'adattatore di \textup{\textsc{Buyer}}.}
\label{tab:bspl_physical_example}

% --- Prima riga ---
\begin{subtable}[t]{0.26\textwidth}
\centering
\caption*{$R(\text{init})$}
\begin{tabular}{l}
\toprule
\textbf{\texttt{ID}} \\
\midrule
\texttt{\#100} \\
\texttt{\#101} \\
\texttt{\#102} \\
\bottomrule
\end{tabular}
\end{subtable}
\hfill
\begin{subtable}[t]{0.34\textwidth}
\centering
\caption*{$R(\text{rfq})$}
\begin{tabular}{ll}
\toprule
\textbf{\texttt{ID}} & \textbf{\texttt{item}} \\
\midrule
\texttt{\#100} & \texttt{"Mouse"} \\
\texttt{\#101} & \texttt{"Libro"} \\
\texttt{\#102} & \texttt{"Tastiera"} \\
\bottomrule
\end{tabular}
\end{subtable}
\hfill
\begin{subtable}[t]{0.34\textwidth}
\centering
\caption*{$R(\text{quote})$}
\begin{tabular}{llr}
\toprule
\textbf{\texttt{ID}} & \textbf{\texttt{item}} & \textbf{\texttt{price}} \\
\midrule
\texttt{\#100} & \texttt{"Mouse"} & 25\,€ \\
\texttt{\#101} & \texttt{"Libro"} & 20\,€ \\
\bottomrule
\end{tabular}
\end{subtable}

\vspace{1.5em} % Spaziatura verticale tra le due righe

% --- Seconda riga ---
\begin{subtable}[t]{0.48\textwidth}
\centering
\caption*{$R(\text{payM})$}
\begin{tabular}{lcr}
\toprule
\textbf{\texttt{ID}} & \textbf{\texttt{price}} & \textbf{\texttt{paid}} \\
\midrule
\texttt{\#100} & 25\,€ & \texttt{true} \\
\bottomrule
\end{tabular}
\end{subtable}
\hfill
\begin{subtable}[t]{0.48\textwidth}
\centering
\caption*{$R(\text{receiptM})$}
\begin{tabular}{lccr}
\toprule
\textbf{\texttt{ID}} & \textbf{\texttt{price}} & \textbf{\texttt{paid}} & \textbf{\texttt{receipt}} \\
\midrule
\texttt{\#100} & 25\,€ & \texttt{true} & \texttt{"REC-112"} \\
\bottomrule
\end{tabular}
\end{subtable}

\end{table}

Confrontando la \zcref{tab:bspl_physical_example} con la vista logica della \zcref{tab:bspl_relational_example}, si nota
che ciascuna relazione locale contiene unicamente tuple popolate con valori definiti per i
messaggi effettivamente scambiati. 
Ad esempio, la chiave \texttt{\#102} non è presente nelle relazioni $R(quote)$ e $R(receiptM)$ poiché i relativi messaggi non 
hanno ancora avuto luogo.
Quando poi l'agente riceverà l'offerta per l'istanza \texttt{\#102} (ad esempio a 45\,€), l'adattatore
eseguirà un semplice inserimento della tupla \texttt{(\#102, "Tastiera", 45\,€)} nella tabella $R(quote)$, senza 
dover effettuare operazioni di modifica in-place sulle altre tabelle.

\subsubsection{Operatività dell'adattatore}\label{bspl:operativita-adattatore}

Quando l'agente invia o riceve un messaggio di schema $m$, l'adattatore deve aggiungere la tupla $t$ 
corrispondente in $R(m)$, per aggiornare lo stato di conoscenza locale.
Tale aggiunta viene effettuata sempre sfruttando una medesima procedura di \textbf{inserimento locale},
la quale, innanzitutto, verifica che i parametri chiave
siano definiti e che non vi siano conflitti di consistenza: relativamente ai parametri condivisi
da $R(m)$ con altre relazioni locali $R(n)$, $t$ deve presentare valori identici a quelli delle tuple
già registrate in $R(n)$ che hanno la sua stessa chiave.
Se queste prime verifiche hanno esito positivo, la procedura controlla che la tupla $t$ non sia già presente
in $R(m)$, ed in caso affermativo la scarta come duplicato senza sollevare
errori né inoltrare notifiche ridondanti alla logica di business.
Diversamente, se la tupla $t$ non era stata già precedentemente inserita, l'adattatore la aggiunge a $R(m)$.

Di seguito si illustrano nel dettaglio le modalità con cui l'adattatore gestisce la ricezione e l'emissione dei messaggi. 

\paragraph{Gestione della ricezione}
A livello di semantica formale, come discusso nel \zcref{bspl:viabilita}, la ricezione 
di un messaggio è incondizionata: un agente non necessita di prerequisiti 
informativi pregressi per accogliere nuovi dati. Tuttavia, in un'architettura 
distribuita ed asincrona concreta, i canali di comunicazione possono essere inaffidabili, 
comportare ritardi arbitrari, recapiti fuori ordine o duplicazioni di pacchetti.

Per questa ragione, l'adattatore, all'arrivo 
di un messaggio $t$ relativo allo schema $m$ dal canale di trasporto, non lo riversa 
direttamente nello stato locale, ma ne verifica la conformità allo schema e in seguito
tenta la procedura di inserimento locale in $R(m)$, durante la quale vengono effettuati dei
controlli ulteriori.
Infine, se l'inserimento ha successo, e dunque il messaggio è valido e non
è un duplicato, l'adattatore notifica la logica di business dell'agente, rendendo 
disponibili i nuovi valori appresi.

\paragraph{Gestione dell'emissione}
Quando invece la logica di business decide di inviare un messaggio di schema $m$, passa all'adattatore la chiave dell'interazione di riferimento
unitamente ai valori calcolati per i relativi parametri \adorn{out}.
Quest'ultimo procede innanzitutto interrogando le relazioni locali per verificare che tutti i parametri \adorn{in}
di $m$ siano già noti per la specifica chiave e che nessun parametro \adorn{out}
o \adorn{nil} risulti già vincolato. 
Se entrambe le verifiche hanno successo, l'adattatore esegue l'inserimento locale della tupla $t$ in $R(m)$,
e trasmette il messaggio sul canale di trasporto. 
In caso contrario, respinge l'invio sollevando un'eccezione locale ed evitando la trasmissione di comunicazioni illegali.

\subsection{Un esempio di protocollo}\label{bspl:esempio-finale}

A completamento della trattazione su BSPL, viene ora presentato un protocollo di riferimento 
denominato \emph{PurchaseWithDelivery}, tratto dalla letteratura su LoST \autocite{singh2011lost}. 

La specifica in questione coinvolge tre ruoli distinti: \textsc{Buyer}, \textsc{Seller} e \textsc{Shipper}. 
I suoi parametri pubblici sono \texttt{ID} (la chiave primaria dell'interazione), \texttt{item} (l'articolo richiesto), 
\texttt{price} (il prezzo quotato) e \texttt{outcome} (l'esito finale della transazione). 
La dichiarazione include inoltre due parametri locali: \texttt{address}, che contiene le informazioni di spedizione, 
e \texttt{response}, utilizzato per gestire la mutua esclusione tra accettazione e rifiuto.

Il protocollo modella una transazione commerciale a tre parti che include una fase iniziale di negoziazione, dove \textsc{Buyer}
interroga \textsc{Seller}, ottenendo il prezzo per un determinato prodotto, e una fase di scelta, dove \textsc{Buyer} comunica a 
\textsc{Seller} la sua volontà di acquistare o meno il prodotto. Infine, se \textsc{Buyer} sceglie di comprare il prodotto
ha luogo anche una fase finale di consegna che coinvolge la terza parte \textsc{Shipper}.

La \zcref{lst:bspl_purchasewithdelivery} riporta la dichiarazione del protocollo in BSPL.

\begin{bspllisting}[caption={Dichiarazione del protocollo PurchaseWithDelivery in BSPL.}, label={lst:bspl_purchasewithdelivery}]
PurchaseWithDelivery {
    role Buyer, Seller, Shipper
    parameter out ID key, out item, out price, out outcome

    Buyer -> Seller: rfq[out ID, out item]
    Seller -> Buyer: quote[in ID, in item, out price]
    Buyer -> Seller: accept[in ID, in item, in price, out address, out response]
    Buyer -> Seller: reject[in ID, in item, in price, out outcome, out response]

    Seller -> Shipper: ship[in ID, in item, in address]
    Shipper -> Buyer: deliver[in ID, in item, in address, out outcome]
}
\end{bspllisting}

\subsubsection{Scelta locale e mutua esclusione}\label{bspl:scelta-locale}

Uno degli aspetti più rilevanti di questo protocollo risiede nella gestione del bivio decisionale tra accettazione 
e rifiuto dell'offerta, che esemplifica il pattern di \textbf{scelta locale} di BSPL \autocite{singh2011bspl}.

Nei protocolli tradizionali basati sul controllo del flusso, la gestione di decisioni alternative distribuite tra 
partecipanti distinti rischia di generare situazioni in cui nodi diversi 
intraprendono percorsi incompatibili all'insaputa l'uno dell'altro.

In BSPL, una scelta tra opzioni mutuamente esclusive viene espressa definendo dei messaggi alternativi 
che condividono i medesimi parametri \adorn{in} e almeno un parametro \adorn{out}, ed affidandone l'emissione allo stesso ruolo.
Nel protocollo \emph{PurchaseWithDelivery}, sia \texttt{accept} che \texttt{reject} richiedono come parametri \adorn{in} 
\texttt{ID}, \texttt{item} e \texttt{price}, e come parametro \adorn{out} \texttt{response}.

Poiché entrambi i messaggi hanno come mittente \textsc{Buyer}, l'invio di uno dei due vincola istantaneamente il 
valore di \texttt{response} nello stato di conoscenza locale dell'acquirente. 
In base alle regole di viabilità, la presenza dell'assegnamento per \texttt{response} rende impossibile
un successivo invio di un messaggio \texttt{accept} o \texttt{reject}.
In questo modo, la mutua esclusione viene garantita localmente, senza richiedere
meccanismi di sincronizzazione o negoziazione con gli altri partecipanti.

\subsubsection{History Vector}\label{bspl:history-vector}

Un'istanza di esecuzione di un protocollo distribuito definito con BSPL è formalmente 
modellata attraverso il concetto di \textbf{history vector} \autocite{singh2011lost}.

Dato un insieme di ruoli $\mathcal{R} = \{x_1, \dots, x_n\}$, un \emph{history vector} $\vec{H}$ è 
definito come un vettore:
\begin{equation}
    \vec{H} = [\mathcal{H}_{x_1}, \dots, \mathcal{H}_{x_n}]
\end{equation}
dove ciascun elemento $\mathcal{H}_{x}$ è la storia locale del ruolo $x$, ossia la sequenza ordinata di messaggi inviati e ricevuti da $x$ nel corso di una specifica esecuzione.

Un \emph{history vector} è valido se rispetta due vincoli fondamentali.
In primo luogo, se un messaggio compare come ricevuto nella storia locale
di un destinatario, deve essere già stato inviato dal rispettivo mittente.
In secondo luogo, se un messaggio compare come inviato all'interno di una 
storia locale, deve risultare viabile rispetto alla conoscenza accumulata
dall'agente fino a quel momento.

La \zcref{fig:esempio_history_vec} illustra un esempio di \emph{history vector} per il protocollo \emph{PurchaseWithDelivery},
relativo a una possibile esecuzione in cui \textsc{Buyer} accetta l'offerta e riceve la merce ordinata.

\begin{figure}[htbp]
    \centering
    \includegraphics[width=0.7\textwidth]{diagrams/esempio_history_vector.pdf}
    \caption{Esempio di esecuzione del protocollo PurchaseWithDelivery tramite history vector.}
    \label{fig:esempio_history_vec}
\end{figure}

L'esecuzione di \emph{PurchaseWithDelivery} mostrata nel diagramma si svolge nei seguenti passi:
\begin{enumerate}
    \item \textsc{Buyer} invia \texttt{rfq}, generando i valori di \texttt{ID} e \texttt{item}.
    \item \textsc{Seller} riceve \texttt{rfq} e risponde inviando \texttt{quote} con il parametro \texttt{price}.
    \item \textsc{Buyer} riceve \texttt{quote} e invia \texttt{accept}, producendo \texttt{address} e \texttt{response}.
    \item \textsc{Seller} riceve \texttt{accept} e inoltra \texttt{ship} al corriere \textsc{Shipper}, trasmettendo \texttt{ID}, \texttt{item} e \texttt{address}.
    \item \textsc{Shipper} riceve \texttt{ship} e conclude l'interazione inviando \texttt{deliver} a \textsc{Buyer}, producendo il parametro \texttt{outcome}.
\end{enumerate}

Come si può notare, ciascun partecipante avanza nella propria storia locale unicamente in base ai dati che riceve o genera, 
senza richiedere una sequenza temporale globale condivisa.